\MFQLead{估计与正则化答案}
\begin{enumerate}[leftmargin=*]
  \item 条件数与实现交叉检查支持数值稳定；抽样和误差协议支持统计识别；经济机制、可交易定义与外部证据支持解释。
  \item 高相关特征可互相替代，轻微样本扰动会改变被保留者。应报告重采样选择频率而非只报一次非零系数。
  \item 梯度为 $X^{\mathsf T}(X\beta-y)+\lambda\beta$，令其为零得到 $(X^{\mathsf T}X+\lambda I)\beta=X^{\mathsf T}y$。
  \item 一维子问题的解是对 $x_j^{\mathsf T}r_{-j}$ 软阈值后除以 $x_j^{\mathsf T}x_j$。
  \item 在相同截距和损失缩放下，透明与库系数应落在声明容差内；否则先检查目标契约。
  \item 零列令坐标更新分母为零，应直接拒绝；近共线列应由奇异值和条件数报告。
  \item 该研究材料把最终期用于选择惩罚参数，因此所谓测试表现有选择偏差；需重新划分未触碰测试期。
  \item WDI 数据不是资产收益，样本极小且无交易机制或因果识别；只可演示外部数据责任链。
\end{enumerate}

\section*{独立 oracle 核对}
透明 ridge 系数为 $0.465574$，透明 lasso 系数为 $0.470000$。
零范数特征列必须在迭代前被拒绝。

\section*{逐题推导与核对}
\begin{enumerate}[leftmargin=*]
  \item 数值稳定先看设计矩阵奇异值、条件数和两种分解的残差；统计识别再看抽样协议、标准误和可辨识参数；经济解释要说明变量如何进入决策、能否交易以及外部机制。一个小残差只能支持第一层，不能自动推出后两层。
  \item lasso 选中是一次样本和一次惩罚下的事件。对高相关列做 bootstrap 或滚动重抽样，记录每列被选中的频率、符号和共同进入次数；若频率分散，结论应称“预测组合稳定”而不是“某个变量解释稳定”。
  \item ridge 目标为 $Q(\beta)=\frac12\|y-X\beta\|_2^2+\frac\lambda2\|\beta\|_2^2$。展开梯度 $\nabla Q=-X^T(y-X\beta)+\lambda\beta$，令零得 $(X^TX+\lambda I)\beta=X^Ty$；$\lambda>0$ 使矩阵在满秩不足时也通常可逆，但会引入收缩偏差。
  \item 固定其他坐标，令 $r_{-j}=y-\sum_{k\ne j}x_k\beta_k$，一维目标为 $\frac12\|r_{-j}-x_j\beta_j\|^2+\lambda|\beta_j|$。若 $z=x_j^Tr_{-j}$、$a=x_j^Tx_j$，分三种情形：$z>\lambda$ 时 $\beta_j=(z-\lambda)/a$；$z<-\lambda$ 时 $\beta_j=(z+\lambda)/a$；$|z|\le\lambda$ 时 $\beta_j=0$，合并即 $S(z,\lambda)/a$。
  \item 惩罚路径先冻结标准化、截距和目标，再扫描一组 $\lambda$。表中同时列每个 $\lambda$ 的 ridge/lasso 系数、非零数、训练/验证损失和条件数；透明实现与库实现只在同一目标、同一缩放下比较，否则差异可能来自契约而非算法。
  \item 全零列的 $x_j^Tx_j=0$，坐标更新没有分母，应在迭代前以 schema/方差门禁拒绝或明确删除；近共线列则报告最小奇异值和条件数，并比较解对微小扰动的变化。不能让求解器静默返回一个看似有限的系数。
  \item 若最终测试期被用来选择 $\lambda$，测试误差条件在选择事件上，已不是无偏的样本外估计。修复是把选择移到内层训练窗，冻结后只访问一次未触碰测试；若测试已被反复查看，应重新采集确认窗并披露旧测试已消耗。
  \item WDI 截面数据的观测对象、频率和缺失机制与股票收益不同，没有资产池、交易时点、成本或因果处理定义。它可以演示数据下载、版本和缺失责任链，但不能支持股票 alpha、风险价格或因果结论；报告必须把演示外推边界写出来。
\end{enumerate}
